\chapter{Tips, Workflow, and Troubleshooting}

\section{Recommended workflow}

For any new dataset, follow this order:

\begin{enumerate}[leftmargin=1.5em]
  \item \textbf{Load} your files (\texttt{.dat} or \texttt{.mat}). Use stride subsampling if the dataset is large.
  \item \textbf{Set} acquisition type and $f_s$ in the popup dialog or the info ribbon.
  \item \textbf{Transform}: rotate to correct tilt, shift the origin to a physical reference.
  \item \textbf{Inspect} the mean field. Does the flow look correct? Right direction, sensible magnitudes?
  \item \textbf{Mean \& Convergence}: plot the mean velocity profile through your region of interest, then check at a representative point that first and second moments are converged. Skip $\langle u'u'u' \rangle$ if your dataset is small.
  \item \textbf{Reynolds stresses}: identify turbulent regions and shear layers.
  \item \textbf{TKE budget}: understand energy production and transport.
  \item \textbf{Correlations}: measure integral length and time scales.
  \item \textbf{POD}: identify dominant coherent structures.
  \item \textbf{DMD} (TR only): extract frequency-specific coherent modes.
  \item \textbf{Vortex identification}: locate and quantify vortex cores.
  \item \textbf{Export} results. Use the multi-case workflow to overlay against benchmark data when relevant.
\end{enumerate}

\section{Performance tips}

\begin{itemize}[leftmargin=1.5em]
  \item POD and DMD computation time scales as $N_t^2$. For 5000+ snapshots use stride subsampling to $\sim$1000--2000 samples first.
  \item Velocity data is loaded as 32-bit floats (float32), halving RAM usage compared to the 64-bit default.
  \item Close unused analysis windows to free RAM.
  \item Use \uipth{Reload Last Dataset} to undo a bad transform without restarting.
  \item For spatial spectra over a region, the FFT tab is faster than running multiple Welch lines.
  \item The $\Gamma_1/\Gamma_2$ computation involves a Python loop over all grid points and can be slow for large grids. Start with a subsampled dataset or a small neighborhood size $S$ while exploring.
  \item For very large \texttt{.mat} files, save with \texttt{-v7.3} in MATLAB to enable HDF5 chunked reading and real per-snapshot progress.
\end{itemize}

\section{Exporting figures for publication}

\begin{itemize}[leftmargin=1.5em]
  \item \textbf{PDF / SVG}: all text is embedded as editable text objects. Open in Adobe Illustrator or Inkscape to change fonts, sizes, and labels to match journal requirements.
  \item \textbf{PNG}: 300 DPI by default. Suitable for Word, PowerPoint, and reports.
  \item Enable \uipth{Clean export} and/or \uipth{Hide colorbar} for plain field images without annotations.
\end{itemize}

\section{Known limitations}

\begin{itemize}[leftmargin=1.5em]
  \item 3D (tomographic) PIV data is not supported.
  \item Pressure and viscous dissipation are not computed; the TKE budget residual absorbs these terms.
  \item Taylor hypothesis (for converting temporal to spatial scales) is only valid at low turbulence intensity ($u'/\langle U\rangle \ll 1$).
  \item Anisotropy analysis requires stereo data (2D3C) and will be unavailable for 2D2C datasets.
  \item DMD requires time-resolved data and a minimum of 50 snapshots.
  \item The $\Gamma_1/\Gamma_2$ computation is slower than gradient-based methods for large grids due to the neighborhood summation loop.
  \item All exported results are time-averaged unless otherwise noted. Per-snapshot export will be added in a future version.
  \item \texttt{.mat} files must contain multiple snapshots. Single-snapshot files are rejected before any data is loaded.
  \item For classic \texttt{.mat} files (not v7.3), the snapshot subset is applied after the full read because \texttt{scipy.io.loadmat} is a single blocking call. Convert to v7.3 in MATLAB for efficient strided loading on very large files.
\end{itemize}

\section{Common issues}

\begin{table}[htbp]
\centering
\caption{Common issues and solutions.}
\begin{tabular}{p{5.2cm}p{8.8cm}}
\toprule
\textbf{Issue} & \textbf{Solution} \\
\midrule
Files do not load (\texttt{.dat}) & Confirm files are Tecplot ASCII format from DaVis. Binary \texttt{.dat} files are not supported. \\
\addlinespace
\texttt{.mat} load shows ``Insufficient Snapshots'' immediately & The file contains only one snapshot. \software{uPrime} requires multiple snapshots. Re-export from MATLAB with the time axis included as the third dimension of \texttt{u}, \texttt{v}, etc. \\
\addlinespace
\texttt{.mat} variable dialog cannot find \texttt{u} & The velocity variable name is not in the recognized list. Use the manual dropdowns in the dialog to map your variables to $x$, $y$, $u$, $v$, etc. \\
\addlinespace
\texttt{.mat} dx/dy shown as 0.000 mm & Coordinate arrays are degenerate (all values equal along one axis), or the meshgrid orientation could not be inferred. Check that \texttt{x} and \texttt{y} contain real coordinate values. \\
\addlinespace
Anisotropy module is greyed out & Data is 2D2C. This module requires stereo (2D3C) data with a $W$ component. \\
\addlinespace
DMD button is greyed out & Set Acquisition Type to TR in the info ribbon and ensure at least 50 snapshots are loaded. \\
\addlinespace
Temporal tabs are greyed out & Set Acquisition Type to Time-Resolved and enter $f_s$ before opening the analysis window. \\
\addlinespace
Top-bar amber non-convergence label is shown & Dataset has fewer than 500 snapshots. Open Mean \& Convergence and verify convergence at a point in your region of interest. \\
\addlinespace
PDF save fails: permission denied & The file is open in Adobe Acrobat. Close it in Acrobat first, then save again. \\
\addlinespace
Rotation gives a blank result & Mirror the dataset \emph{before} rotating, not after. \\
\addlinespace
Lumley triangle points outside boundary & Insufficient snapshots or masked/invalid vectors near the sampling line. Avoid masked regions when drawing the analysis line. \\
\addlinespace
DMD spectrum dominated by high-$|\sigma|$ modes & Apply the growth rate filter ($|\sigma| <$ threshold) to focus on near-neutral physically meaningful modes. \\
\addlinespace
$\Gamma$ computation is very slow & Reduce the neighborhood half-width $S$, or subsample the dataset before computing. \\
\addlinespace
Imported case appears at wrong scale & The imported CSV uses a different $U_m$ or different units than the current session. Check the case's header and re-export consistently. \\
\bottomrule
\end{tabular}
\end{table}

% ============================= CHAPTER 15 ==============================
